\documentclass[10pt]{article}
\usepackage[utf8]{inputenc}
\usepackage[T1]{fontenc}
\usepackage{amsmath}
\usepackage{amsfonts}
\usepackage{amssymb}
\usepackage[version=4]{mhchem}
\usepackage{stmaryrd}
\usepackage{graphicx}
\usepackage[export]{adjustbox}
\graphicspath{ {./images/} }

\begin{document}
Таким образом, отображение $d$ есть эпиморфпзм мультипликативных полугрупп $M_{n}(K)$ и $K$.

Пусть $m \leqslant n, A=\left\|a_{i} j\right\|$ есть $(m \times n)$-матрица, $B=\left\|b_{i}\right\|$ есть $(n \times m)$-матрица над $K$, а $C=A B$. Тогда верна фо р м ула Бин е - К оши:

\[
\begin{aligned}
& \operatorname{det} C=\sum_{1 \leqslant j_{1}<\ldots<j_{m} \leqslant n}\left|\begin{array}{c}
a_{1 j_{1}} \ldots a_{1 j_{m}} \\
. \cdots \\
a_{m j_{1}} \ldots a_{m j m}
\end{array}\right| \times \\
& \left|\begin{array}{c}
b_{j_{1} 1} \ldots b_{j_{1} m} \\
\times \cdot \cdots \cdot \\
b_{j_{m} 1} \cdots b_{j_{m} m}
\end{array}\right|
\end{aligned}
\]

Пусть $A=\left|a_{i j}\right| \in M_{n}(K)$, а $A_{i j}-$ алгебраич. дополнение элемента $a_{i j}$. Тогда верны формулы

\[
\left.\begin{array}{c}
\Sigma_{j=1}^{n} a_{i j} A_{k j}=\delta_{i k} \operatorname{det} A, \\
\Sigma_{i=1}^{n} a_{i j} A_{i k}=\delta_{j k} \operatorname{det} A,
\end{array}\right\}
\]

где $\delta_{i j}-$ символ Кронекера. Для вычислений О. часто используются разложение его по әлементам строки или столбца, т. е. формулы (1), теорема Лапласа (см. А.лгебраическое дополнение.) и преобразования матрицы $A$, ве меняющие 0 . Для матрицы $A$ из $M_{n}(K)$ тогда и только тогда существует обратная матрица $A^{-1}$ в $M_{n}(K)$, когда в $K$ имеется элемент, обратныї әлементу $\operatorname{det} A$. Следовательно, отображение

\[
\mathrm{GL}(n, K) \longrightarrow K^{*}(A \mapsto \operatorname{det} A),
\]

где GL $(n, K)$ - группа всех обратимых матриц в $M_{n}(K)$, т. е. полная линейная группа, а $K^{*}-$ группа обратимых элементов $K$, есть эпиморфизм этих групп.

Квадратная матрица над полем обратима тогда и только тогда, когда ее О. отличен от нуля. $n$-мерные векторы $a_{1}, \ldots, a_{n}$ над полем $F$ линейно зависимы тогда и только тогда, когда

\[
D\left(a_{1}, \ldots, a_{n}\right)=0 .
\]

О. матрицы $A$ порядка $n>1$ над полем равен 1 тогда и только тогда, когда $A$ есть произведение элементарных матриц вида

\[
t_{i j}(\lambda)=E_{n}+\lambda e_{i j}
\]

где $i \neq j$, а $e_{i j}$ - матрица, единственный ненулевої элемент к-рой равен 1 и расположен на позиции $(i, j)$.

Теория О. возникла в связи с задачей решения систем линейных уравнений:

\[
\left.\begin{array}{c}
a_{11} x_{1}+\ldots+a_{1 n} x_{n}=b_{1}, \\
a_{n 1} x_{1}+\ldots+a_{n n} x_{n}=b_{n},
\end{array}\right\}
\]

где $a_{i j}, b_{j}$ - элементы нек-рого поля $F$. Если $\operatorname{det} A \neq 0$, где $A=\left\|a_{i j}\right\|$ - матрица системы (2), то эта система имеет единственное решение, вычисляемое по формулам Крамера (см. К рамера правило). В случае, когда система (2) задана над кольцом $K$ и $\operatorname{det} A$ обратим в $K$, система также имеет единственное решение, определяемое теми же формулами Крамера.

Теория О. построена также и для матриц над некоммутативным ассоциативным телом. О. матрицы над телом $k$ (о п р е д е л и т е л ь Д ь ё д он н е) вводится следующим образом. Тело $k$ рассматривается как полугруппа и строится ее коммутативный гомоморфіный образ $\bar{k} . k-$ группа $k *$ с внешне присоединенным нулем 0 , а в качестве $\bar{k}$ берется также группа $\bar{k}^{*} \mathrm{c}$ внешне присоединенным нулем $\overline{0}$, где $\bar{k} *$ - факторгруппа группы $k$ * по коммутанту. Эпиморфизм $k \rightarrow \bar{k}$, $\lambda \rightarrow \bar{\lambda}$, задается канонич. эпиморфизмом групп $k^{*} \rightarrow \overline{k^{*}}$ и условием $0 \rightarrow \overline{0}$. Очевидно, $\overline{1}-$ единица полугруппы $\bar{k}$.

Теория О. над телом основана ва следующей теореме. Существует единственное отображение

\[
\delta: M_{n}(k) \mapsto \bar{k},
\]

удовлетворяющее следующим трем аксиомам:

I) если матрица $B$ получена из матрицы $A$ умножением слева одной строки на $\lambda \in k$, то $\delta(B)=\bar{\lambda} \delta(A)$;

II) если $B$ получена из $A$ заменой строки $a_{i}$ строкой $a_{i}+a_{j}$, где $i \neq j$, то $\delta(B)=\delta(A)$;

III) $\delta\left(E_{n}\right)=\overline{1}$.

Элемент $\delta(A)$ наз. о п р е д е л и т е л е м матрицы $A$ и обозначается $\operatorname{det} A$. Для коммутативного тела ак сиомы I), II), III) совпадают с условиями $1_{\mathrm{a}}$ ), 2), 3) соответственно, и, следовательно, в этом случае получаются обычные О. над полем. Если $A=\operatorname{diag}\left[a_{11}, \ldots\right.$, $\left.a_{n n}\right]$, то $\operatorname{det} A=\overline{a_{11} \ldots a_{n n}}$; таким образом, отображение $\delta: M_{n}(k) \rightarrow \bar{k}$ сюръективно. Матрица $A$ из $M_{n}(k)$ обратима тогда и только тогда, когда $\operatorname{det} A \neq 0$. Справедливо равенство $\operatorname{det} A B=\operatorname{det} A \cdot \operatorname{det} B$. Как и в коммутативном случае, $\operatorname{det} A$ не изменится, если строку $a_{i}$ матрицы $A$ заменить строкой $a_{i}+\lambda a_{j}$, где $i \neq j, \lambda \in k$. При $n>1 \operatorname{det} A=1$ тогда и только тогда, когда $A$ произведение әлементарных матриц вида $t_{i j}(\lambda)=E_{n}+\lambda e_{i j}$, $i \neq j, \lambda \in k$. Если $a \neq 0$, то

\[
\left|\begin{array}{ll}
a & b \\
c & d
\end{array}\right|=\overline{a d-a c a^{-1} b},\left|\begin{array}{ll}
0 & b \\
c & d
\end{array}\right|=-\overline{c b} .
\]

В отличие от коммутативного случая, $\operatorname{det} A^{T}$ может и не совпадать с $\operatorname{det} A$. Напр., для матрицы $A=\left\|\begin{array}{cc}i & j \\ k-1\end{array}\right\|$ над телом кбитернионов $\operatorname{det} A=-\overline{2 i}$, а $\operatorname{det} A^{T}=\overline{0}$.

Бесконечные О., то есть О. бесконечных матриц, определяются как предел, к к-рому стремится О. конечной подматрицы при бесконечном возрастании ее порядка. Если әтот предел существует, то О. наз. сходящимся, в противном случае - расходящимся.

Понятие «O.» восходит к Г. Лейбницу (G. Leibnitz, 1678); первая публикация принацлежит Г. Крамеру (G. Cramer, 1750). Теория О. создана трудами А. Вандермонда (A. Vandermonde), П. Лапласа (P. Laplace), О. Коши (А. Cauchy) и К. Якоби (С. Јасobi). Термин «O.» встречается впервые у К. Гаусса (C. Gauss, 1801). Современное обозначение введено А. Кэли (А. Cayley, $1841)$.

Лит.: [1] К у р о ші А. Г., Курс высшей алгебры, 11 изд.,

\includegraphics[max width=\textwidth, center]{2023_07_08_d1363a0d7b0850c00b26g-1}
алгебра и многомерная геометрия, М., $1970 ;$ [4]' Т ы ш к ев и ч Р. И., Ф е д е н к о А. С., Линейная алгебра и аналитическая геометрия, 2 изд., Минск, $1976 ;[5]$ А р т и н Э., Геометри-
ческая алгебра, пгер. с англ., М., $1969 ;$ [6] Б у р б а к и Н., ческая алгебра, пер. с англ., М., 1969 ; [6] Б у р б а к и Н., алгебра, пер. с франц., М., 1962; [7] К а г а н В. Ф., Основания теории определителей, Одесса, 1922. Д. А. Супруненко.

ОПРЕДЕЛЯ ЮШАЯ СИСТЕМА ОКРЕСТНОСТЕЙ мно жест в а $A$ в тополо ги ч е ском пр ос т р а н с т в е $X$-любое семейство $\xi$ подмножеств пространства $X$, подчиненное следующим двум условиям: а) для каждого $O \in \xi$ найдется открытое множество $V$ пространства $X$ такое, что $O \supset V \supset A$, б) каково бы ни было открытое в $X$ множество $W$, содержащее $A$, найдется элемент $U$ семейства $\xi$, содержащийся в $W$. Иногда дополнительно предполагают, что все элементы семейства $\xi$ открытые множества. Определяющей системой окрестностей точки $x \in X$ в топологич. пространстве $X$ наз. О. с. о. в $X$ одноточечного множества $\{x\}$. Лum.: [1] А р х анге ль с ки й А. В., П о ном аре в В. И., Основы общей топологии в задачах и упражнениях,
М., 1974.


\end{document}